\section*{Hoofdstuk \ref{ch:g12:reaction-rates} --- Reactiesnelheden}

\begin{solution}{exo:g12:reaction-rates:1}
Temperatuur; contactoppervlak; concentratie; temperatuur.
\end{solution}

\begin{solution}{exo:g12:reaction-rates:2}
\qty{13.9}{min} (blauw) en \qty{6.9}{min} (rood). Na twee
\omterm{def:g12:reaction-rates:half-life}{halveringstijden}: $10.0 / 4 = \qty{2.5}{mmol/L}$.
\end{solution}

\begin{solution}{exo:g12:reaction-rates:3}
Helling $(5.0 - 1.0)/10 = \qty{0.40}{mmol.L^{-1}.min^{-1}}$: dat is de
\omterm{def:g12:reaction-rates:rate}{vormingssnelheid} bij \qty{5}{min}.
\end{solution}

\begin{solution}{exo:g12:reaction-rates:4}
Afkoelen en verdunnen maken de reactie zo langzaam dat de samenstelling
tijdens de \omterm{def:g11:titration:titration}{titratie} niet meer verandert. De gemeten samenstelling is die
van het \omterm{def:g6:pure-substances-and-mixtures:mixture}{mengsel} op het moment dat de reactie werd gestopt, bij het
uitgieten.
\end{solution}

\begin{solution}{exo:g12:reaction-rates:5}
\qty{30}{min} is 3 \omterm{def:g12:reaction-rates:half-life}{halveringstijden}: er blijft $1/8$ over. \qty{1}{h} is
6 \omterm{def:g12:reaction-rates:half-life}{halveringstijden}: $1/64$, ongeveer \qty{1.6}{\percent}.
\end{solution}

\begin{solution}{exo:g12:reaction-rates:6}
$\ln[\ce{A}]$: 2.079, 1.733, 1.386, 1.040, 0.693. Het daalt met 0.346 per
\qty{5}{min}: een rechte lijn, eerste orde. $k = 0.346/5 =
\qty{0.0693}{min^{-1}}$ en $t_{1/2} = \ln 2 / k = \qty{10.0}{min}$ (wat je
ook direct ziet: 8.00, 4.00, 2.00 om de \qty{10}{min}).
\end{solution}

\begin{solution}{exo:g12:reaction-rates:7}
$k = 0.693 / 25 = \qty{0.0277}{s^{-1}}$.
\end{solution}

\begin{solution}{exo:g12:reaction-rates:8}
$0.100\,e^{-0.20} = \qty{0.0819}{mol/L}$; $0.100\,e^{-0.60} =
\qty{0.0549}{mol/L}$; $0.100\,e^{-2.0} = \qty{0.0135}{mol/L}$.
\end{solution}

\begin{solution}{exo:g12:reaction-rates:9}
Bij \qty{13.9}{min} is $[\ce{A}] = \qty{5.0}{mmol/L}$; de helling van de
raaklijn is ongeveer $-0.25$, en $-k[\ce{A}] = -0.050 \times 5.0 =
\qty{-0.25}{mmol.L^{-1}.min^{-1}}$: de helft van de beginwaarde.
\end{solution}

\begin{solution}{exo:g12:reaction-rates:10}
Jood is het enige gekleurde deeltje: volgens de wet van Lambert-Beer is
$A = \varepsilon \ell [\ce{I2}]$. De \omterm{def:g12:reaction-rates:rate}{vormingssnelheid} is de helling van
de raaklijn aan de kromme $[\ce{I2}](t)$, dat is de helling van de
raaklijn aan $A(t)$ gedeeld door $\varepsilon \ell$.
\end{solution}

\begin{solution}{exo:g12:reaction-rates:11}
De beginsnelheid $k[\ce{A}]_0$ is voor de tweede proef twee keer zo
groot; de \omterm{def:g12:reaction-rates:half-life}{halveringstijden} zijn gelijk, $\ln 2 / k$, onafhankelijk van
$[\ce{A}]_0$.
\end{solution}

\begin{solution}{exo:g12:reaction-rates:12}
$t = \ln(1/f)/k$. In \omterm{def:g12:reaction-rates:half-life}{halveringstijden}: $\ln(1/f)/\ln 2$. Voor
\qty{1}{\percent}: $\ln 100 / \ln 2 = 6.6$; voor \qty{0.1}{\percent}:
$\ln 1000 / \ln 2 =
10.0$.
\end{solution}

\begin{solution}{exo:g12:reaction-rates:13}
\qty{90}{\percent} verbruikt, $f = 0.10$: $t = \ln 10 / k$, dat is
\qty{46}{min} bij \qty{20}{\celsius} en \qty{23}{min} bij
\qty{30}{\celsius}. De \omterm{def:g12:reaction-rates:first-order}{snelheidsconstante} verdubbelt over
\qty{10}{\celsius}: de regel klopt hier.
\end{solution}

\begin{solution}{exo:g12:reaction-rates:14}
De temperatuur daalt tot dicht bij \qty{0}{\celsius}, en de concentraties
worden door 10 gedeeld. Bij een snelheid die evenredig is met het product
van twee concentraties deelt de \omterm{def:g10:concentration-and-dilution:dilution}{verdunning} alleen haar al door
$10 \times 10 = 100$.
\end{solution}

\begin{solution}{exo:g12:reaction-rates:15}
$v(t) = -\dfrac{d}{dt}\big([\ce{A}]_0 e^{-kt}\big) = k[\ce{A}]_0 e^{-kt}$.
Beginsnelheid $k[\ce{A}]_0$; de snelheid is gehalveerd als
$e^{-kt} = 1/2$, dat is na één \omterm{def:g12:reaction-rates:half-life}{halveringstijd}.
\end{solution}

\begin{solution}{pb:g12:reaction-rates:1}
\textbf{1.} Door de wet van Lambert-Beer, omdat de kleurstof het enige
deeltje is dat bij deze golflengte absorbeert; de oplossing moet verdund
genoeg zijn ($A$ onder ongeveer 1).

\textbf{2.} Het absorbeert geen zichtbaar licht: het draagt niets bij
aan $A$.

\textbf{3.} Zodat de concentratie ervan nauwelijks verandert: de snelheid
hangt dan alleen van de kleurstof af.

\textbf{4.} $A = 0$ (de blanco).

\textbf{5.} $A$ daalt in \qty{6}{min} van 0.800 naar 0.402: ongeveer
\qty{6}{min}.

\textbf{6.} Van 0.402 bij \qty{6}{min} naar 0.199 bij \qty{12}{min}:
opnieuw gehalveerd in \qty{6}{min}.

\textbf{7.} $-0.223$, $-0.448$, $-0.691$, $-0.911$, $-1.155$, $-1.366$,
$-1.614$, $-1.945$, $-2.551$, $-3.079$.

\textbf{8.} De punten liggen op een rechte lijn: de reactie is van de
eerste orde in de kleurstof.

\textbf{9.} $(-2.551 + 0.223)/20 = \qty{-0.116}{min^{-1}}$.

\textbf{10.} $k = \qty{0.116}{min^{-1}}$.

\textbf{11.} $0.693 / 0.116 = \qty{6.0}{min}$, zoals in de tabel
afgelezen.

\textbf{12.} $k A_0 = 0.116 \times 0.800 = \qty{0.093}{min^{-1}}$ in
absorbantie-eenheden per minuut.

\textbf{13.} Dezelfde, \qty{6.0}{min}: een \omterm{def:g12:reaction-rates:half-life}{halveringstijd} bij de eerste
orde hangt niet af van de beginwaarde.

\textbf{14.} $0.800\, e^{-0.116 \times 30} = 0.025$.

\textbf{15.} $t = \ln 100 / k = 4.61 / 0.116 = \qty{40}{min}$.

\textbf{16.} $40 / 6.0 = 6.6$ \omterm{def:g12:reaction-rates:half-life}{halveringstijden}.

\textbf{17.} $t = \ln(0.800/0.010)/k = 4.38 / 0.116 = \qty{38}{min}$.

\textbf{18.} Nog steeds 6.6 \omterm{def:g12:reaction-rates:half-life}{halveringstijden}, nu
$6.6 \times 3.0 = \qty{20}{min}$.

\textbf{19.} Ongeveer 6.6 \omterm{def:g12:reaction-rates:half-life}{halveringstijden}, dat is hier ongeveer
\qty{40}{min}.
\end{solution}
